\documentclass[11pt,reqno]{amsart}
\usepackage[T1]{fontenc}\usepackage[utf8]{inputenc}\usepackage{lmodern}
\usepackage[a4paper,margin=24mm]{geometry}
\usepackage{amsmath,amssymb,mathtools,microtype,xurl}
\usepackage[hidelinks]{hyperref}
\newtheorem{theorem}{Theorem}\newtheorem{corollary}[theorem]{Corollary}
\title{A Complement--Norm Cutoff for Exterior Erd\H{o}s--Straus States}
\author{The Clankers}\date{19 September 2026}
\begin{document}
\begin{abstract}
For an original exterior state, retain its residual $R$, cofactor $D$, and
complementary square divisor $v$. We prove the strict bound
$(p+w)^2>4(w+1)^2(w-1)$ with $w=\min(R,D,v)$. This gives a complete three-chart
atlas with outer cutoff $(p/2)^{2/3}+O(p^{1/3})$, including its original
availability condition and denominator inverse. A hard-prime example separates
this channel from the preceding middle square-root bound. The argument is a
necessary state reduction; it does not establish universal occupancy.
\end{abstract}\maketitle

Let $p\equiv1\pmod4$, $p>1$, and
\[
 p/4<a<p/2,\quad\gcd(p,a)=1,\quad R=4a-p,\quad
 u\mid a^2,\quad R\mid4u+1.
\]
Gcd normalization gives
\[
 g=\gcd(a,u),\quad(h,r,s)=(g^2/u,u/g,a/g),\quad
 \kappa=(pr+s)/R,
\]
\[
 a=hrs,\quad u=hr^2,\quad \gcd(r,s)=1,\qquad
 \frac4p=\frac1a+\frac1{hs\kappa}+\frac1{phr\kappa}.
\]
The ordered inverse is $u=a^2/(Ry-pa)$. These are the classical Type-I
coordinates [1]. Put
\[
 v=a^2/u=hs^2,\qquad
 D=(4u+1)/R=(r+\kappa)/s.
\]
Then
\[
 pD+1=4hr\kappa,\quad (p+R)^2+4v=4vRD,\quad p+R<2vD.       \tag{1}
\]
The last inequality follows from $\kappa>r$, itself a consequence of $R<p$.
Both $R,D$ are $3\pmod4$. The case $h=1$ would make $-1$ a square modulo $R$;
hence $h>1$. Since $h$ is coprime to $R,D$, the integer $v$ equals neither.
This does not assert that $v$ is coprime to $D$.

\begin{theorem}
Every stated source satisfies
\[
 \boxed{(p+w)^2>4(w+1)^2(w-1),\qquad w=\min(R,D,v).}         \tag{2}
\]
Let $C(p)$ be the largest integer satisfying (2). Then $C(p)<p$ and
$C(p)=(p/2)^{2/3}+O(p^{1/3})$.
\end{theorem}
\begin{proof}
Set $t=\min(R,D)$. If $R=t$, (1) gives
$(p+t)^2\ge4v(t^2-1)$. If $D=t\le R$, subtract to obtain
\[
 (p+t)^2-4v(t^2-1)=(R-t)(4vt-2p-R-t)\ge0;
\]
the second factor is positive by (1). Thus
$p\ge2\sqrt{v(t^2-1)}-t$, with equality in the first step only if $R=D$.

If $v>t$, integrality gives $v\ge t+1$ and therefore
$p\ge2(t+1)\sqrt{t-1}-t$. Equality would make $t-1$ an integer square,
whereas $t\equiv3\pmod4$. It is impossible.

If $v<t$, put $w=v$. The preceding expression increases in $t$, so
$p\ge2w\sqrt{w+2}-w-1$. This is strictly greater than
$2(w+1)\sqrt{w-1}-w$: the difference between
$w\sqrt{w+2}$ and $(w+1)\sqrt{w-1}$ is
\[
 \frac{w^2+w+1}{w\sqrt{w+2}+(w+1)\sqrt{w-1}}>\frac12.
\]
Squaring the positive sides proves (2). The bounding function is increasing
for $w>1$; at $w=p$ its required inequality fails since $p^3-p-1>0$.
Its expansion is $2w^{3/2}-w+O(\sqrt w)$, proving the asymptotic statement.
The integer cutoff is evaluated by comparisons in (2), not by that expansion.
\end{proof}

For $v=\prod\ell^{e_\ell}$ define
$K(v)=\prod\ell^{\lceil e_\ell/2\rceil}$. The equivalence
$v\mid a^2\iff K(v)\mid a$ is an exact prime-valuation identity.

\begin{theorem}[Complete prime atlas]
For prime $p\equiv1\pmod4$, all original first-half E states occur in the
following disjoint charts, with $C=C(p)$.

\textup{(R)} Use the original square-divisor source at
$R\equiv3\pmod4$, $3\le R\le C$, and retain $R\le D$, $R<v$.

\textup{(D)} At $D\equiv3\pmod4$, $3\le D\le C$, put $B=(pD+1)/4$.
Take $w\mid B^2$, $w<B$, $D\mid w+B$; gcd-normalize
\[
 h=\gcd(B,w)^2/w,\quad r=w/\gcd(B,w),\quad\kappa=B/\gcd(B,w),
\]
then $s=(r+\kappa)/D$, $a=hrs$, $u=w$. Retain $D<R$, $D<v$.

\textup{(v)} For $2\le v\le C$, take
\[
 \boxed{R\mid p^2+4v,\quad v<R<p,\quad R\equiv3\pmod4,
 \quad R\equiv-p\pmod{4K(v)}.}                          \tag{3}
\]
Set $a=(p+R)/4$, $u=a^2/v$, $D=(4u+1)/R$ and retain $v<D$.
Gcd normalization gives the original state.
\end{theorem}
\begin{proof}
For (D), normalization gives $r<\kappa$, $4hr\kappa=pD+1$ and
$D\mid r+\kappa$. Then $R\kappa=pr+s>0$ and
\[
 D(p-2a)=2hr(\kappa-r)-1>0.
\]
Thus the original first-half inequalities hold. The inverse reads
$B=hr\kappa$, $w=hr^2$ with its retained orientation.

In (v), $v<p$ and primality imply $\gcd(R,4v)=1$. The last congruence in
(3) gives the actual availability $v\mid a^2$. Now
\[
 4v(4u+1)=(p+R)^2+4v\equiv0\pmod R
\]
gives the original gate after cancellation of a unit. Conversely (1) gives
$R\mid p^2+4v$ for each original state. All inverses are literal.
The smallest parameter is at most $C$. The grade $v$ cannot tie $R,D$;
a possible tie $R=D$ is assigned to (R). This proves disjoint completeness.
\end{proof}

\paragraph{Negative controls and scale.}
The hard prime $p=48049\equiv169\pmod{840}$ has the original state
\[
 (a,u,h,r,s,\kappa,R,D,v)
 =(12090,24180,6045,2,1,309,311,311,6045),
\]
with ordered denominators $(12090,1867905,179501934690)$. Both exterior
cofactors exceed the earlier middle cutoff $247$ [2]. At $p=37$, the state
$(a,u,h,r,s,\kappa)=(12,8,2,2,3,7)$ has $(v,D)=(18,3)$; their overlap
cannot be cancelled as though they were coprime.

For every integer $n\ge2$, the exact source
\[
 h=4n^2-2,\quad r=n,\quad s=1,\quad R=D=4n^2-1,
 \quad\kappa=4n^2-n-1
\]
has $p=16n^3-4n^2-8n+1$, $a=hn$, $u=hn^2$ and $\gcd(p,a)=1$.
Its least parameter is $h\sim(p/2)^{2/3}$, and in fact $C(p)=h$: the bounding
function at $h+1$ exceeds $p$ because
$(h+2)\sqrt h>h\sqrt{h+2}$. The integer cutoff is attained at the prime
members $p=97,929$. These are legitimate unit integer
sources, with hard residue $289\pmod{840}$ when $n=24+210j$.
No infinitude of prime values is claimed. The asymptotic sharpness assertion
is in this integer domain.

The complete main replay compares the three charts with original divisor
enumeration; the second checker uses primitive $h,r,s$ enumeration and imports
no first-checker code. The actual finite bound and command outcomes are in the
release receipt. These checks do not imply universal ES positivity. No
historical-priority or independently reviewed-proof claim is made.

\begin{thebibliography}{9}
\bibitem{1} Elsholtz, C., \& Tao, T. (2013). Counting the number of solutions to the
Erd\H{o}s--Straus equation on unit fractions. \emph{Journal of the Australian
Mathematical Society, 94}(1), 50--105. \texttt{arXiv:1107.1010}, \S2.
\bibitem{2} The Clankers. (2026). \emph{A sharp two-chart middle-selector cutoff
and an original-grade obstruction for Erd\H{o}s--Straus}. Preserved Turn 7 source,
sections 1--4. A complete decision atlas, not universal occupancy.
\end{thebibliography}
\end{document}
